% ============================================================================
%  sections/introduction.tex
% ============================================================================
\chapter*{Introduction générale}
\addcontentsline{toc}{chapter}{Introduction générale}

Dans un contexte de transformation numérique des établissements d'enseignement supérieur au Maroc, la gestion des processus d'admission et de présélection constitue un enjeu majeur pour les écoles d'ingénieurs. L'ENSA Béni Mellal accueille chaque année un nombre croissant de candidatures pour ses filières d'ingénieur, tant au niveau Bac+2 qu'au niveau Bac+3.

Le processus actuel de présélection, principalement manuel et basé sur des fichiers Excel, présente plusieurs limites : risque d'erreurs humaines, lenteur du traitement, absence de traçabilité et difficulté de gestion des volumes importants de dossiers. Le service de la Scolarité a exprimé le besoin d'une solution informatique capable d'automatiser et de fiabiliser ce processus.

C'est dans ce contexte que s'inscrit notre projet de stage : la conception et le développement d'une plateforme web intelligente de présélection des candidats. Ce rapport est structuré en quatre chapitres :

\begin{itemize}[leftmargin=2cm]
  \item \textbf{Chapitre 1 --- Contexte du Projet :} présentation de l'organisme d'accueil, problématique et objectifs.
  \item \textbf{Chapitre 2 --- Architecture et Base de données :} architecture Full-Stack, stack technologique et dictionnaire des modèles Django.
  \item \textbf{Chapitre 3 --- Logique Métier et Algorithmes :} moteur de règles, scoring pondéré, pipeline OCR et fuzzy matching.
  \item \textbf{Chapitre 4 --- Implémentation, Asynchronisme et Tests :} interfaces, threading, notifications WebSocket, sécurité et validation.
\end{itemize}
